\documentclass[11pt,a4paper]{article}

\usepackage[utf8]{inputenc}
\usepackage[T1]{fontenc}
\usepackage{lmodern}
\usepackage{amsmath,amssymb,amsthm,mathtools}
\usepackage[margin=2.5cm]{geometry}
\usepackage{hyperref}
\usepackage{cleveref}
\usepackage{enumitem}
\usepackage{booktabs}
\usepackage{xcolor}
\usepackage{fancyhdr}

\newcommand{\dd}{\mathrm{d}}
\newcommand{\Weylsq}{C_{\mu\nu\rho\sigma}C^{\mu\nu\rho\sigma}}
\newcommand{\Rsq}{R^2}

\newtheorem{theorem}{Theorem}[section]
\newtheorem{proposition}[theorem]{Proposition}
\newtheorem{remark}[theorem]{Remark}
\newtheorem{definition}[theorem]{Definition}

\pagestyle{fancy}
\fancyhf{}
\fancyhead[L]{\small\textsc{SCT Theory --- PPN-1}}
\fancyhead[R]{\small\thepage}
\renewcommand{\headrulewidth}{0.4pt}

\hypersetup{
  colorlinks=true,
  linkcolor=blue!60!black,
  citecolor=green!40!black,
  urlcolor=blue!60!black,
  pdftitle={PPN-1 in SCT Theory: Linear Static Yukawa Sector},
  pdfauthor={David Alfyorov},
}

\title{%
  \textbf{PPN-1 in SCT Theory}\\[6pt]
  \large Linear Static Sector from the NT-4a Propagator
}
\author{David Alfyorov}
\date{March 2026}

\begin{document}
\maketitle


\begin{center}
\small\textit{Credits: Aliaksandr Samatyia contributed research-assistance and workflow support.}
\end{center}

\begin{abstract}
This document derives the part of the PPN-1 program that is currently under
control in SCT Theory. The input is not the full nonlinear post-Newtonian
system but the linearized NT-4a propagator obtained from the spectral
curvature-squared action. We first write the exact linear static potentials
as Fourier--Bessel integrals in terms of the propagator denominators
$\Pi_{\mathrm{TT}}(z)$ and $\Pi_{\mathrm{s}}(z,\xi)$. We then pass to the
local Yukawa approximation, which is the regime actually implemented in the
current code and sufficient for first-pass Solar-System and short-range
gravity estimates.

Within that restricted but honest scope, the derivation fixes the radial
profile of the linear PPN parameter $\gamma(r) = \Psi(r)/\Phi(r)$ and the
associated lower bounds on the SCT scale $\Lambda$ from Cassini and
short-range inverse-square-law tests. The document also states clearly what
is \emph{not} derived here: the nonlinear PPN parameter $\beta$, the
preferred-frame parameters $\alpha_i$, the preferred-location parameter
$\xi_{\mathrm{PPN}}$, and the conservation-law parameters $\zeta_i$.
Those quantities require the $\mathcal{O}(h^2)$ field equations and
therefore remain open at the present stage.
\end{abstract}

\tableofcontents
\newpage

\section{Scope and Conventions}

\subsection{What This Phase Does}

The purpose of PPN-1 is to connect the verified spectral form factors and the
NT-4a propagator with observable weak-field quantities. In the present repair
cycle, the phase is narrowed to the piece that is actually supported by the
existing derivation chain:
\begin{enumerate}[label=\arabic*., leftmargin=2em]
  \item the exact \emph{linear static} response of the metric,
  \item the \emph{local Yukawa approximation} to that response,
  \item the induced profile of $\gamma(r)$,
  \item the resulting first bounds on $\Lambda$.
\end{enumerate}

The following are \emph{not} claimed in this document:
\begin{enumerate}[label=\alph*)]
  \item a full ten-parameter PPN completion,
  \item a derivation of $\beta$,
  \item a proof that the preferred-frame or conservation parameters vanish,
  \item an analytic evaluation of the exact nonlocal Fourier integrals at all
    radii.
\end{enumerate}

\subsection{Metric and Sign Conventions}

We use the Lorentzian weak-field form
\begin{equation}
  g_{00} = -1 + 2\Phi(r), \qquad
  g_{ij} = \bigl(1 + 2\Psi(r)\bigr)\delta_{ij},
\end{equation}
with signature $(-,+,+,+)$. In the static limit, the physical potential
energy of a non-relativistic test body is
\begin{equation}
  V(r) = -m_{\mathrm{test}}\,\Phi(r).
\end{equation}
Thus the Newtonian potential corresponds to
\begin{equation}
  \Phi_{\mathrm{N}}(r) = \Psi_{\mathrm{N}}(r) = \frac{GM}{r},
  \qquad
  V_{\mathrm{N}}(r) = -\frac{GMm_{\mathrm{test}}}{r}.
\end{equation}

\subsection{Dimensional Conventions}

Natural units $c=\hbar=1$ are used in the derivation. For phenomenology we
restore the conversion
\begin{equation}
  1\,\mathrm{m} \simeq 5.06773\times 10^{6}\,\mathrm{eV}^{-1}.
\end{equation}
In particular,
\begin{equation}
  1\,\mathrm{AU}
  \simeq 1.495978707\times 10^{11}\,\mathrm{m}
  \simeq 7.58\times 10^{17}\,\mathrm{eV}^{-1},
\end{equation}
and
\begin{equation}
  50\,\mu\mathrm{m}
  \simeq 2.53\times 10^{2}\,\mathrm{eV}^{-1}.
\end{equation}

\section{Input from Phase 3 and NT-4a}

The starting point is the verified spectral action in the Weyl basis,
\begin{equation}
  \Gamma^{(1)}
  = \frac{1}{16\pi^2}\int \dd^4x \,\sqrt{g}\,
  \left[
    \alpha_C \hat F_1\!\left(\frac{\Box}{\Lambda^2}\right)\Weylsq
    + \alpha_R(\xi)\hat F_2\!\left(\frac{\Box}{\Lambda^2}\right)\Rsq
  \right],
\end{equation}
with
\begin{equation}
  \alpha_C = \frac{13}{120},
  \qquad
  \alpha_R(\xi) = 2\left(\xi - \frac{1}{6}\right)^2,
\end{equation}
and normalized form factors
\begin{equation}
  \hat F_i(z) \coloneqq \frac{F_i(z)}{F_i(0)},
  \qquad
  \hat F_i(0)=1.
\end{equation}

The NT-4a propagator denominators are
\begin{equation}
  \Pi_{\mathrm{TT}}(z) = 1 + c_2 z \hat F_1(z),
  \qquad
  c_2 = \frac{13}{60},
\end{equation}
\begin{equation}
  \Pi_{\mathrm{s}}(z,\xi)
  = 1 + \bigl(3c_1 + c_2\bigr) z \hat F_2(z,\xi)
  = 1 + 6\left(\xi-\frac{1}{6}\right)^2 z \hat F_2(z,\xi).
\end{equation}
Two structural facts matter immediately:
\begin{enumerate}[nosep]
  \item $\Pi_{\mathrm{TT}}(0)=1$ and $\Pi_{\mathrm{s}}(0,\xi)=1$,
  so the infrared limit is Einsteinian.
  \item At conformal coupling $\xi = 1/6$, one has
  $\Pi_{\mathrm{s}}(z,1/6)=1$, so the scalar channel decouples from the
  linearized dynamics.
\end{enumerate}

\section{Exact Linear Static Potentials}

\subsection{Fourier Representation}

For a point source
\begin{equation}
  T_{00} = M \delta^{(3)}(\vec x),
  \qquad
  T_{0i}=T_{ij}=0,
\end{equation}
the linearized NT-4a solution can be written in terms of the exact
spin-2 and spin-0 denominators:
\begin{equation}\label{eq:Phi-exact}
  \Phi(r) = \frac{GM}{r}\,\frac{2}{\pi}
  \int_0^\infty \frac{\sin(kr)}{k}
  \left[
    \frac{4}{3\,\Pi_{\mathrm{TT}}(k^2/\Lambda^2)}
    - \frac{1}{3\,\Pi_{\mathrm{s}}(k^2/\Lambda^2,\xi)}
  \right]\dd k,
\end{equation}
\begin{equation}\label{eq:Psi-exact}
  \Psi(r) = \frac{GM}{r}\,\frac{2}{\pi}
  \int_0^\infty \frac{\sin(kr)}{k}
  \left[
    \frac{2}{3\,\Pi_{\mathrm{TT}}(k^2/\Lambda^2)}
    + \frac{1}{3\,\Pi_{\mathrm{s}}(k^2/\Lambda^2,\xi)}
  \right]\dd k.
\end{equation}
These expressions are the exact linearized SCT answer for the static metric.

\begin{remark}
At this stage the nonlocal theory is still present in full. The only
approximation made so far is the restriction to the linearized static sector.
No local pole approximation has yet been taken. In the current normalized
NT-4a implementation, the TT denominator also develops a positive-real-axis
zero near $z \approx 2.41484$. This means the exact static kernel cannot be
advertised as ghost-free on the basis of the present data alone.
\end{remark}

\subsection{What Would Be Needed for Full PPN}

The standard PPN expansion requires:
\begin{enumerate}[label=\arabic*., leftmargin=2em]
  \item $g_{00}$ through $\mathcal{O}(v^4)$,
  \item $g_{0i}$ through $\mathcal{O}(v^3)$,
  \item $g_{ij}$ through $\mathcal{O}(v^2)$,
  \item the nonlinear superposition structure of the field equations.
\end{enumerate}
Equations \eqref{eq:Phi-exact}--\eqref{eq:Psi-exact} determine only the
linear static part of the PPN data. This is enough for $\gamma(r)$ but not
for $\beta$ or the full preferred-frame sector.

\section{Local Yukawa Approximation}

\subsection{Small-\texorpdfstring{$z$}{z} Expansion}

For $k^2 \ll \Lambda^2$, we use
\begin{equation}
  \hat F_1(z) = 1 + \mathcal{O}(z), \qquad
  \hat F_2(z,\xi) = 1 + \mathcal{O}(z),
\end{equation}
so the denominators become
\begin{equation}
  \Pi_{\mathrm{TT}}(z) = 1 + c_2 z + \mathcal{O}(z^2),
\end{equation}
\begin{equation}
  \Pi_{\mathrm{s}}(z,\xi)
  = 1 + \bigl(3c_1+c_2\bigr)z + \mathcal{O}(z^2).
\end{equation}
This motivates the effective masses
\begin{equation}
  m_2^2 = \frac{\Lambda^2}{c_2}
  = \frac{60}{13}\Lambda^2,
\end{equation}
\begin{equation}
  m_0^2 = \frac{\Lambda^2}{3c_1+c_2}
  = \frac{\Lambda^2}{6(\xi-1/6)^2},
\end{equation}
with $m_0 = \infty$ when $\xi = 1/6$.

\subsection{Local Potentials}

Replacing the exact denominators by their Yukawa pole approximants gives
\begin{equation}\label{eq:Phi-local}
  \Phi_{\mathrm{loc}}(r)
  = \frac{GM}{r}
  \left[
    1 - \frac{4}{3}e^{-m_2 r} + \frac{1}{3}e^{-m_0 r}
  \right],
\end{equation}
\begin{equation}\label{eq:Psi-local}
  \Psi_{\mathrm{loc}}(r)
  = \frac{GM}{r}
  \left[
    1 - \frac{2}{3}e^{-m_2 r} - \frac{1}{3}e^{-m_0 r}
  \right].
\end{equation}
These formulas are the actual content of the repaired implementation in
\path{analysis/scripts/nt4a_newtonian.py} and
\path{analysis/scripts/ppn1_parameters.py}.

\begin{proposition}[Short-distance structure of the local approximation]
For $\xi \neq 1/6$, the scalar Yukawa term cancels the leading $1/r$
singularity:
\begin{equation}
  \Phi_{\mathrm{loc}}(r)
  = GM\left(\frac{4}{3}m_2 - \frac{1}{3}m_0\right) + \mathcal{O}(r),
  \qquad r \to 0.
\end{equation}
At exact conformal coupling, the cancellation is absent and
\begin{equation}
  \Phi_{\mathrm{loc}}(r)\big|_{\xi=1/6}
  \sim -\frac{GM}{3r},
  \qquad r \to 0.
\end{equation}
\end{proposition}

\begin{proof}
Expand $e^{-mr} = 1 - mr + \mathcal{O}(r^2)$ in
\cref{eq:Phi-local,eq:Psi-local}. For $\xi \neq 1/6$ one has both Yukawa
terms, and the constant coefficients satisfy $1 - 4/3 + 1/3 = 0$.
At exact conformal coupling the scalar term is absent from the start,
leaving $1 - 4/3 = -1/3$ and hence a residual $1/r$ behavior.
\end{proof}

\begin{remark}
The conformal and short-distance limits do not commute. Taking
$\xi \to 1/6$ at fixed $r>0$ removes the scalar Yukawa tail, whereas
taking $r \to 0$ first at non-conformal $\xi$ shows the cancellation of the
Newtonian singularity. This point must be kept explicit in any later UV
interpretation.
\end{remark}

\section{The Linear PPN Parameter \texorpdfstring{$\gamma(r)$}{gamma(r)}}

\subsection{Derivation}

The linear PPN parameter is defined by the ratio of spatial and temporal
metric potentials:
\begin{equation}
  \gamma(r) \coloneqq \frac{\Psi(r)}{\Phi(r)}.
\end{equation}
Within the local Yukawa approximation, \cref{eq:Phi-local,eq:Psi-local}
therefore imply
\begin{equation}\label{eq:gamma-local}
  \boxed{
  \gamma_{\mathrm{loc}}(r)
  =
  \frac{
    1 - \frac{2}{3}e^{-m_2 r} - \frac{1}{3}e^{-m_0 r}
  }{
    1 - \frac{4}{3}e^{-m_2 r} + \frac{1}{3}e^{-m_0 r}
  }}.
\end{equation}
This is the only PPN parameter presently derived from the repaired phase.

\subsection{Infrared Limit}

\begin{proposition}[GR recovery]
For $r\Lambda \gg 1$,
\begin{equation}
  \gamma_{\mathrm{loc}}(r) = 1 + \mathcal{O}(e^{-m_2 r}, e^{-m_0 r}).
\end{equation}
\end{proposition}

\begin{proof}
Both exponentials vanish at large radius, so numerator and denominator of
\cref{eq:gamma-local} approach unity.
\end{proof}

\subsection{Short-Distance Values}

For non-conformal $\xi$, the local $r \to 0$ limit is
\begin{equation}
  \gamma_{\mathrm{loc}}(0)
  =
  \frac{2m_2 + m_0}{4m_2 - m_0}.
\end{equation}
At exact conformal coupling the scalar mode is absent and
\begin{equation}
  \gamma_{\mathrm{loc}}(0)\big|_{\xi=1/6} = -1.
\end{equation}
These short-distance values are not directly constrained by Solar-System
observations, but they matter for the internal consistency of the local
approximation and for matching to short-range laboratory tests.

\section{Phenomenological Bounds on \texorpdfstring{$\Lambda$}{Lambda}}

\subsection{Cassini Bound}

The Cassini time-delay measurement requires
\begin{equation}
  |\gamma - 1| < 2.3\times 10^{-5}
\end{equation}
at roughly $r \simeq 1\,\mathrm{AU}$. In the Yukawa regime with
$m_0 \geq m_2$, the dominant correction comes from the spin-2 term:
\begin{equation}
  |\gamma_{\mathrm{loc}} - 1|
  \approx \frac{2}{3}e^{-m_2 r}.
\end{equation}
Demanding compatibility with Cassini gives
\begin{equation}
  e^{-m_2 r_{\mathrm{AU}}}
  \lesssim 3.45\times 10^{-5},
\end{equation}
hence
\begin{equation}
  m_2 r_{\mathrm{AU}} \gtrsim 10.28.
\end{equation}
Using $r_{\mathrm{AU}} \simeq 7.58\times 10^{17}\,\mathrm{eV}^{-1}$,
we obtain
\begin{equation}
  m_2 \gtrsim 1.36\times 10^{-17}\,\mathrm{eV},
\end{equation}
and therefore
\begin{equation}
  \boxed{
  \Lambda \gtrsim 6.2\times 10^{-18}\,\mathrm{eV}
  \qquad \text{(Cassini, local Yukawa estimate)}.}
\end{equation}

\subsection{Short-Range Inverse-Square-Law Tests}

Torsion-balance searches for Yukawa corrections at
$\lambda \lesssim 50\,\mu\mathrm{m}$ imply
\begin{equation}
  m_2 \gtrsim \frac{1}{50\,\mu\mathrm{m}}
  \simeq 3.95\times 10^{-3}\,\mathrm{eV}.
\end{equation}
Since $m_2 = \Lambda \sqrt{60/13}$, this gives
\begin{equation}
  \boxed{
  \Lambda \gtrsim 1.84\times 10^{-3}\,\mathrm{eV}
  \qquad \text{(short-range gravity, local Yukawa estimate)}.}
\end{equation}
This laboratory bound is much stronger than the Cassini one.

\begin{remark}
These are bounds on the \emph{local Yukawa approximation}, not on the exact
entire-function propagator. They are therefore safe as first constraints but
should later be recomputed from the full nonlocal static kernel.
\end{remark}

\section{What Remains Open}

\subsection{The Parameter \texorpdfstring{$\beta$}{beta}}

The nonlinear PPN parameter $\beta$ probes the $\mathcal{O}(h^2)$ sector of
$g_{00}$ and therefore depends on the nonlinear field equations, not merely
on the linear propagator. Since the present phase uses only the linear NT-4a
kernel, no trustworthy derivation of $\beta$ is available yet.

\subsection{Preferred-Frame and Conservation Parameters}

The parameters
\begin{equation}
  \xi_{\mathrm{PPN}}, \qquad \alpha_1,\alpha_2,\alpha_3,
  \qquad \zeta_1,\zeta_2,\zeta_3,\zeta_4
\end{equation}
are not determined by the present calculation. Diffeomorphism invariance is
an important structural input, but it is not by itself a substitute for a
complete PPN reduction of the nonlinear equations. Accordingly, the repaired
code now marks these parameters as \texttt{not\_derived}.

\section{Summary}

The repaired status of PPN-1 can be summarized as follows:
\begin{enumerate}[label=\arabic*., leftmargin=2em]
  \item The exact \emph{linear static} potentials are determined by the NT-4a
  denominators \eqref{eq:Phi-exact}--\eqref{eq:Psi-exact}.
  \item The current implementation uses the \emph{local Yukawa approximation}
  to those exact kernels.
  \item Within that approximation, the linear PPN parameter $\gamma(r)$ is
  given by \cref{eq:gamma-local}.
  \item The current lower bounds are
  \[
    \Lambda \gtrsim 6.2\times 10^{-18}\,\mathrm{eV}
    \quad \text{(Cassini)},
    \qquad
    \Lambda \gtrsim 1.84\times 10^{-3}\,\mathrm{eV}
    \quad \text{(short range)}.
  \]
  \item $\beta$ and the remaining eight PPN parameters are \emph{not yet
  derived}. Any earlier wording that presented them as completed results must
  be regarded as superseded by this repaired document.
\end{enumerate}

\begin{thebibliography}{9}

\bibitem{Will2014}
C. M. Will,
\textit{The Confrontation between General Relativity and Experiment},
Living Rev. Relativity \textbf{17}, 4 (2014).

\bibitem{Stelle1977}
K. S. Stelle,
\textit{Renormalization of Higher-Derivative Quantum Gravity},
Phys. Rev. D \textbf{16}, 953 (1977).

\bibitem{Edholm2016}
J. Edholm, A. S. Koshelev, and A. Mazumdar,
\textit{Universality of the Newtonian Potential in Higher Derivative Gravity},
Phys. Rev. D \textbf{94}, 104033 (2016).

\end{thebibliography}

\end{document}
